% ==================== 1.1 研究背景及意义 ====================
\section{研究背景及意义}
\label{section:1.1}

月球作为距离地球最近的天体，承载着人类深空探测的重要使命。随着航天技术的飞速发展，月面探测已从单纯的科学研究向资源开发、基地建设等方向拓展。月面巡视器作为月面探测的核心装备，其自主行走能力直接决定了探测任务的效率、安全性和可持续性。

\subsection{月面探测的科学价值与工程风险}
\label{section:1.1.1}

月球是地球唯一的天然卫星，蕴含着丰富的科学信息和战略资源。从科学研究角度看，月球保存了太阳系形成早期的地质记录，其表面物质成分、撞击坑分布、火山活动遗迹等为研究地球起源、太阳系演化提供了珍贵样本\cite{ref1}。从资源开发角度看，月球富含氦-3、钛铁矿、稀土元素等战略资源，其中氦-3被认为是未来可控核聚变的理想燃料\cite{ref2}。

然而，月面探测面临严峻的工程风险。月球表面环境极端恶劣：昼夜温差可达300\degree C以上，月尘具有强附着性和磨蚀性，月壤力学特性复杂多变，地形起伏大且存在大量撞击坑和岩石障碍\cite{ref3}。这些因素对巡视器的行走性能提出了极高要求，任何导航决策失误都可能导致巡视器陷入困境甚至任务失败。

\subsection{地-月通讯时延与遥操作局限}
\label{subsec:comm_delay}

地-月通讯存在约2.5秒的往返时延，这一时延对巡视器的实时操控构成了根本性限制。传统的遥操作模式下，地面控制人员发送指令后需要等待反馈信息才能判断执行结果，这种"指令-执行-反馈"的循环模式导致巡视器工作效率极其低下。

以"玉兔号"巡视器为例，其单次移动操作需要经过"图像获取-地面规划-指令上传-执行反馈"等环节，整个过程耗时可达数十分钟\cite{ref4}。更为严峻的是，在通讯中断或信号干扰情况下，地面将完全失去对巡视器的控制能力，巡视器将陷入"盲走"状态，面临极高的安全风险。

因此，提升巡视器的自主行走能力，减少对地面控制的依赖，已成为深空探测技术发展的必然趋势。

\subsection{数字孪生技术的深空探测定位}
\label{subsec:dt_position}

数字孪生（Digital Twin）技术作为一种连接物理世界与虚拟世界的桥梁，为解决上述难题提供了新的技术途径\cite{ref5}。数字孪生通过构建物理实体的虚拟镜像，实现对物理实体的精确映射、状态监测、行为预测和优化控制，其核心价值在于：

（1）\textbf{全要素数字化映射}：将巡视器的几何模型、动力学模型、感知模型、环境模型等要素进行数字化表达，构建与物理实体高度一致的虚拟实体；

（2）\textbf{实时数据驱动}：通过传感器数据实时更新虚拟实体状态，实现虚拟与现实之间的双向同步；

（3）\textbf{智能分析与决策}：在虚拟空间中进行仿真分析、行为预测、策略优化，为物理实体的自主决策提供支持。

将数字孪生技术应用于月面巡视器自主行走，可以在虚拟空间中构建高保真的月面环境和巡视器模型，实现环境感知、路径规划、避障决策等功能的数字化重构，从而在通讯受限条件下支撑巡视器的自主行走能力\cite{ref6}。这一技术路线对于提高巡视器工作效率、保障任务安全、拓展探测范围具有重要的战略意义。
